\section{Introduction}
By interconnecting devices, vehicles, and appliances, 
the Internet of Things (IoT) has transformed sectors ranging from smart cities~\cite{qian2019internet}, which optimize traffic and energy use, to autonomous vehicles~\cite{kang2018blockchain} that aim for safer roads, industrial networks~\cite{singh2020deep} that redefine production processes, e-health solutions~\cite{gadekallu2021blockchain} that prioritize patient care, and smart homes~\cite{park2017comprehensive} that enrich daily living. The massive and sensitive nature of the data generated demands strong security and integrity, leading to the integration of IoT and blockchain~\cite{wang2023blockchain, iot-bc-1, iot-bc-2, 8989788, wang2019survey}. This convergence, with its decentralized, immutable, and traceable features, offers a foundational structure for secure data exchanges~\cite{mathur2023survey}. Sharding, a key strategy in blockchain, addresses vast data needs and ensures scalability~\cite{yu2020survey}. By dividing the network into smaller segments, sharding enables simultaneous transaction processing, lightening the load on nodes and increasing transaction throughput.

As the IoT landscape expands, managing and securing vast networks becomes more intricate. However, while various sharding techniques like random-based~\cite{luu2016secure, kokoris2018omniledger, zamani2018rapidchain, wang2019monoxide}, community-based~\cite{zhang2022community}, and trust-based~\cite{yun2019trust} methods enhance the scalability and distribution of dishonest nodes, they remain vulnerable to advanced threats such as adaptive collusion attacks. In such attacks, dishonest nodes manipulate sharding systems by modifying their behavior to focus on specific shards, potentially initiating the $1\%$ attack~\cite{aponte202151}, which significantly compromises the blockchain's security, trustworthiness, and immutability.

Deep Reinforcement Learning (DRL) has become instrumental in the confluence of IoT and blockchain sharding due to its proficiency in handling multiple sharding variables~\cite{yun2020dqn, yang2022sharded,liu2019performance, liu2018blockchain, qiu2019service, qiu2018blockchain}. Amidst the surging blockchain transaction data, DRL provides real-time adaptability, fostering dynamic shard adjustments that enhance resource use and overall system efficiency. However, many DRL-focused studies underestimate the challenges of strategic collusion attacks, where dishonest nodes might disguise their identity or increase cross-shard transactions to hide their intent.

This study investigates the expansion and fortification of blockchain sharding by introducing an innovative framework that utilizes trust tables and DRL. The introduced trust table advances the sharding process by incorporating a multi-dimensional approach to gather feedback within the voting mechanism. This multi-faceted feedback includes direct feedback, indirect feedback, and historical behaviors during the voting process, all contributing to a robust defense mechanism against $1\%$ attack risks. Concurrently, the DRL-driven sharding framework is designed to dynamically allocate shards in real-time, which streamlines node synchronization and minimizes cross-shard transactions (CSTs). This approach also ensures a balanced distribution of dishonest nodes, aiming to decrease the necessity for frequent resharding. The study suggests security thresholds for this DRL approach to enhance training and fortify the sharded blockchain's security.

The primary contributions of this study can be summarized as follows:

\begin{itemize}

\item We propose a secure, scalable, sharded blockchain system architecture for IoT scenarios. Our proposed system effectively mitigates the risks associated with $1\%$ attacks and reduces CST, striking a balance between security and scalability.

\item  We design a novel blockchain sharding system, \textsc{TbDd}, incorporating trust tables to distinguish between dishonest and honest nodes. We consider multiple perspectives for generating a trust table to identify node properties, such as distributed voting, consensus, and node historical behavior analysis.

\item We leverage the DRL framework to monitor the blockchain sharding process continually. By dynamically balancing scalability and security, we achieve high throughput while maintaining the number of dishonest nodes below the security threshold in each shard, ensuring a robust and secure decentralized system.
\end{itemize}

Comparative experiments were conducted simulating strategic collusion attacks across different sharding methods. The outcomes highlight that the \textsc{TbDd} system excels in minimizing CST and ensuring a balanced shard trust distribution. Additionally, the \textsc{TbDd} framework, in the absence of collusion attacks and within tolerable dishonest node levels, offers about a $10\%$ throughput advantage over random-based sharding and a $13\%$ enhancement against the trust-based approach.


The organization of the remaining paper is as follows: Section~\ref{section: related work} reviews relevant literature. Section~\ref{section: system model} introduces the proposed system model and problem definition. Section~\ref{section: TBDD} comprehensively presents the proposed \textsc{TbDd} system. Experiments and assessments are conducted in Section~\ref{section: experiment}. The paper is concluded in Section~\ref{section: conclusion}.
